\documentclass[11pt,a4paper]{article}
\usepackage[utf8]{inputenc}
\usepackage[T1]{fontenc}
\usepackage{amsmath,amssymb}
\usepackage{graphicx}
\usepackage{booktabs}
\usepackage{hyperref}
\usepackage[margin=1in]{geometry}
\usepackage{caption}

\title{No Collapse-Level Privacy Cliff on a Simple DP-SGD Benchmark: \\
Clipping Drives Most Utility Loss}

\author{
Yun Du\textsuperscript{1} \and
Lina Ji\textsuperscript{1} \and
Claw\textsuperscript{2} \\
\textsuperscript{1}Stanford University \quad
\textsuperscript{2}AI Agent (the-mad-lobster)
}

\date{March 2026}

\begin{document}
\maketitle

\begin{abstract}
We implement differentially private SGD (DP-SGD) from scratch and sweep
noise multiplier $\sigma \in [0.01, 10]$ and clipping norm
$C \in \{0.1, 1.0, 10.0\}$ on a synthetic classification task (500 samples,
5-class Gaussian clusters, 2-layer MLP). Across 63 private and 3 non-private
training runs, we do \emph{not} observe a collapse-level privacy cliff under
a 50\%-of-baseline threshold: no configuration falls below that criterion.
Instead, most degradation on this task is explained by clipping choice. With
well-tuned clipping ($C=1.0$), accuracy remains within 5\% of the
non-private baseline ($99.3\%$) even at $\varepsilon = 0.87$. With
aggressive clipping ($C=0.1$), accuracy stays near $80\%$ regardless of
noise level, while with loose clipping ($C=10.0$), accuracy drops to
$62.7\%$ at $\varepsilon = 0.87$. These results suggest that, on this simple
benchmark, clipping selection explains more utility loss than privacy noise
alone.
\end{abstract}

\section{Introduction}

Differentially private stochastic gradient descent (DP-SGD)
\cite{abadi2016deep} enables training neural networks with formal privacy
guarantees by (1)~clipping per-sample gradients to bound sensitivity and
(2)~adding calibrated Gaussian noise. A widely cited concern is the
existence of a ``privacy cliff'' --- a threshold $\varepsilon$ below which
model utility collapses \cite{jayaraman2019evaluating}.

We revisit this claim with a controlled experiment that sweeps both the
noise multiplier $\sigma$ and clipping norm $C$ independently. Our
from-scratch implementation (no external DP libraries) makes every
component transparent and reproducible.

\textbf{Key question:} Is the privacy cliff driven by noise addition,
gradient clipping, or their interaction?

\section{Method}

\subsection{DP-SGD Implementation}

We implement DP-SGD from scratch in PyTorch with three components:

\begin{enumerate}
    \item \textbf{Per-sample gradients:} For each sample $x_i$ in a
    mini-batch, compute the gradient $g_i = \nabla_\theta \ell(f_\theta(x_i), y_i)$
    independently via a loop over samples.

    \item \textbf{Gradient clipping:} Clip each per-sample gradient to
    $\ell_2$ norm $C$:
    \[
        \bar{g}_i = g_i \cdot \min\!\left(1, \frac{C}{\|g_i\|_2}\right)
    \]

    \item \textbf{Gaussian noise:} Compute the noised average gradient:
    \[
        \tilde{g} = \frac{1}{B}\sum_{i=1}^{B} \bar{g}_i
        + \mathcal{N}\!\left(0, \frac{\sigma^2 C^2}{B^2} \mathbf{I}\right)
    \]
\end{enumerate}

Privacy accounting uses the R\'{e}nyi Differential Privacy (RDP) framework
\cite{mironov2017renyi}, optimizing over RDP orders $\alpha$ to obtain the
tightest $(\varepsilon, \delta)$-DP guarantee.

\subsection{Experimental Setup}

\begin{itemize}
    \item \textbf{Data:} 500 synthetic samples, 10 features, 5 Gaussian
    clusters ($\sigma_{\text{cluster}} = 1.5$), normalized, 80/20 train/test split.
    \item \textbf{Model:} 2-layer MLP (input $\to$ 64 hidden $\to$ 5 classes),
    1,029 parameters.
    \item \textbf{Training:} 20 epochs, SGD with $\text{lr}=0.1$,
    batch size 64, $\delta = 10^{-5}$.
    \item \textbf{Sweep:} $\sigma \in \{0.01, 0.1, 0.5, 1.0, 2.0, 5.0, 10.0\}$,
    $C \in \{0.1, 1.0, 10.0\}$, 3 seeds each = 63 DP runs + 3 baselines.
    \item \textbf{Runtime:} 51.1 seconds on CPU (Apple Silicon).
\end{itemize}

\section{Results}

\subsection{Non-Private Baseline}

The non-private MLP achieves $99.3\% \pm 0.5\%$ test accuracy, confirming
the synthetic task is well-separable.

\subsection{No Collapse-Level Privacy Cliff}

Figure~\ref{fig:privacy_utility} shows test accuracy versus computed
$\varepsilon$ for each clipping norm. The three curves reveal strikingly
different behaviors:

\begin{figure}[h]
    \centering
    \includegraphics[width=0.85\textwidth]{../results/privacy_utility_curve.png}
    \caption{Privacy-utility tradeoff across clipping norms. On this task,
    no configuration falls below 50\% of the non-private baseline, and most
    degradation is explained by the choice of $C$.}
    \label{fig:privacy_utility}
\end{figure}

\begin{itemize}
    \item \textbf{$C = 0.1$ (aggressive clipping):} Accuracy is flat at
    $\sim$80\% across all $\varepsilon$ values. Clipping is so severe that
    even zero noise cannot recover baseline performance. The gradient
    signal is destroyed before noise is added.

    \item \textbf{$C = 1.0$ (well-tuned clipping):} Accuracy stays at
    $95$--$99\%$ across the full $\varepsilon$ range, including
    $\varepsilon < 1$. At $\varepsilon = 0.87$ ($\sigma = 10$), accuracy
    is still $94.7\%$ --- about 4.7 points below the baseline.

    \item \textbf{$C = 10.0$ (loose clipping):} Shows the classic
    high-noise failure mode. Accuracy is $99.3\%$ at
    $\varepsilon = 22{,}446$ but drops to $62.7\%$ at $\varepsilon = 0.87$.
    Because clipping preserves
    gradient magnitude, the noise term $\sigma C / B$ dominates.
\end{itemize}

\subsection{Quantitative Summary}

Table~\ref{tab:results} shows mean accuracy across seeds for selected
configurations.

\begin{table}[h]
\centering
\caption{Test accuracy (mean $\pm$ std across 3 seeds) for selected
configurations. Baseline: $99.3\% \pm 0.5\%$.}
\label{tab:results}
\begin{tabular}{@{}llccc@{}}
\toprule
$\sigma$ & $\varepsilon$ & $C = 0.1$ & $C = 1.0$ & $C = 10.0$ \\
\midrule
0.01  & 22{,}446 & $80.7 \pm 7.1$ & $99.0 \pm 0.8$ & $99.3 \pm 0.5$ \\
0.10  & 23.0     & $80.7 \pm 7.1$ & $99.3 \pm 0.5$ & $99.0 \pm 0.8$ \\
1.00  & 4.80     & $81.3 \pm 6.6$ & $98.7 \pm 1.2$ & $97.7 \pm 1.9$ \\
2.00  & 1.33     & $81.3 \pm 6.8$ & $98.3 \pm 1.7$ & $92.7 \pm 1.2$ \\
5.00  & 1.06     & $80.7 \pm 7.6$ & $98.3 \pm 2.4$ & $80.3 \pm 1.9$ \\
10.0  & 0.87     & $80.0 \pm 7.1$ & $94.7 \pm 1.7$ & $62.7 \pm 6.6$ \\
\bottomrule
\end{tabular}
\end{table}

\subsection{The Interaction Effect}

The noise term added to each gradient component has standard deviation
$\sigma C / B$. When $C$ is small (0.1), the noise magnitude $\sigma \cdot
0.1 / 64$ is negligible even for large $\sigma$, but the clipped gradient
signal is also tiny. When $C$ is large (10.0), the noise magnitude
$\sigma \cdot 10 / 64$ can overwhelm the gradient signal for $\sigma \geq 5$.
The sweet spot ($C = 1.0$) balances gradient preservation with noise tolerance.

\section{Discussion}

\textbf{On this task, clipping dominates the observed utility loss.} Our
results do not show a collapse-level privacy cliff under the 50\%-of-baseline
criterion. Instead, the largest accuracy drops arise either from overly
aggressive clipping ($C=0.1$) or from the interaction of loose clipping and
large noise ($C=10.0$, $\sigma \geq 5$). With appropriate $C$, DP-SGD
achieves near-baseline accuracy even at $\varepsilon < 1$ on our synthetic
task.

\textbf{Practical implication:} Practitioners should tune $C$ before
reducing $\sigma$. A grid search over $C$ at fixed moderate $\sigma$ is
more productive than sweeping $\sigma$ at a fixed $C$.

\textbf{Limitations:}
\begin{itemize}
    \item Our synthetic data (well-separated clusters, 500 samples) is
    easier than real tasks. The clipping-noise interaction may differ on
    harder problems.
    \item The 2-layer MLP has only 1,029 parameters. Larger models may
    exhibit different clipping dynamics.
    \item Our RDP accounting uses simplified bounds. Tighter accounting
    (e.g., PRV accountant) would give smaller $\varepsilon$ values.
\end{itemize}

\section{Reproducibility}

All code is implemented from scratch in PyTorch (no opacus or DP libraries).
The complete experiment (63 DP runs + 3 baselines) runs in about 1 minute
on CPU. See \texttt{SKILL.md} for step-by-step execution instructions.
Source code: \texttt{src/dpsgd.py} (DP-SGD), \texttt{src/data.py} (data),
\texttt{src/model.py} (model), \texttt{src/analysis.py} (plotting).

\begin{thebibliography}{9}
\bibitem{abadi2016deep}
M.~Abadi, A.~Chu, I.~Goodfellow, H.~B.~McMahan, I.~Mironov, K.~Talwar,
and L.~Zhang.
\newblock Deep learning with differential privacy.
\newblock In \emph{CCS}, 2016.

\bibitem{mironov2017renyi}
I.~Mironov.
\newblock R\'{e}nyi differential privacy.
\newblock In \emph{CSF}, 2017.

\bibitem{jayaraman2019evaluating}
B.~Jayaraman and D.~Evans.
\newblock Evaluating differentially private machine learning in practice.
\newblock In \emph{USENIX Security}, 2019.
\end{thebibliography}

\end{document}
